\documentclass[11pt]{article}

% ============================================================================
% intro_cohomology_refactor_v3.tex   (supersedes v1, v2)
%
% v3 CHANGES (per discussion):
%  (a) "Gauge" vocabulary scrubbed from all blocks. Standard cohomological
%      language throughout: redundancy / choice of representative /
%      normalization. The operative phrasing is yours: choose the
%      representative C^f - \delta P whose T-value vanishes.
%  (b) Block 1, par. 2 rebuilt around the referee objection ("why not just
%      the finite S-contour?") and its one-line answer: torsion forces
%      C_S|(1+S) = 0 at every basepoint, and ker(1+S) = im(delta_S), so the
%      raw S-value lies inside its own ambiguity space. The two Kohnen--
%      Zagier relations have different statuses: the S-relation is free,
%      the U-relation is purchased by the T-normalization.
%  (c) Single-P emphasis in par. 1 compressed to one clause.
%  (d) Why T: infinite order, unipotent on V_n (single Jordan block), so
%      delta_T has 1-dim kernel and cokernel -- one scalar obstruction
%      (the residue at the cusp T stabilizes), one scalar of slack.
%
% CONSISTENCY SWEEP (gauge language remaining in the MAIN file, outside
% the replaced spans; convert or keep deliberately, but uniformly):
%   lines 412--413, 469, 589, 600 (Section 1.x prose), and the
%   conclusions block at lines 3060--3079, which mirrors the old intro
%   and should be re-mirrored on the new one.
%
% Blocks:
%   BLOCK 1 -- replaces lines 313--351; delete the %DAM note at 321.
%              Existing line 353 ("The coboundary equation admits a
%              closed-form inverse...") follows directly.
%   BLOCK 2 -- replaces lines 767--774 (incl. %DAM note at 774).
%   BLOCK 3 -- sign fix in eq:bp-cocycle-coboundary (lines 755--760).
%
% \ref/\cite target the MAIN file and render as ??/[?] here by design.
% Open check still open: +/- <-> parity assignment in thm:periods.
% ============================================================================

\usepackage[margin=1in]{geometry}
\usepackage{amsmath, amssymb, amsthm}
\usepackage{mathtools}
\usepackage{xcolor}
\usepackage{hyperref}

\hypersetup{colorlinks=true, linkcolor=blue!60!black, urlcolor=blue!60!black, citecolor=blue!60!black}

% --- macros copied from the main file ---
\newcommand{\HH}{\mathbb{H}}
\newcommand{\ZZ}{\mathbb{Z}}
\newcommand{\QQ}{\mathbb{Q}}
\newcommand{\CC}{\mathbb{C}}
\newcommand{\SLZ}{\mathrm{SL}_2(\ZZ)}
\newcommand{\Dhat}{\widehat{\Delta}}
\DeclareMathOperator{\im}{im}
\DeclareMathOperator{\Hpar}{H^1_{\mathrm{par}}}

% --- presentation-only helpers for this standalone document ---
\newcommand{\blockrule}{\par\noindent\rule{\textwidth}{0.8pt}\par}
\newenvironment{whereitgoes}
  {\par\medskip\noindent\begin{minipage}{\textwidth}\small\itshape\color{black!70}}
  {\end{minipage}\par\medskip}

\title{\bf Refactor of the cohomology passage, v3:\\ standard language,
free-vs-purchased relations}
\author{Editing companion; supersedes v1/v2}
\date{\today}

\begin{document}
\maketitle

\noindent\textbf{How to use this file.}
Block 1 replaces lines 313--351 (delete the answered \texttt{\%DAM} note
at 321; line 353 follows directly).  Block 2 replaces lines 767--774 in
Section~\texttt{sec:cocycle} (delete the \texttt{\%DAM} note at 774).
Block 3 is a one-line sign fix at \texttt{eq:bp-cocycle-coboundary}.
Copy what lies between the \verb|% ==== BEGIN/END PASTE ====| markers in
the source.  \verb|\ref|/\verb|\cite| target the main file and render
here as ``??''/``[?]'' by design.  The source header lists the lines in
the main file where leftover ``gauge'' language should be swept for
consistency (notably the conclusions block at 3060--3079, which mirrors
the old intro).

% ============================================================================
\section*{Block 1 --- replaces lines 313--351 (intro, \S1.1)}
% ============================================================================
\blockrule

% ==== BEGIN PASTE: BLOCK 1 (replaces lines 313--351; delete the %DAM note at line 321) ====

The price of trading one divergent contour for a family of finite ones
is redundancy: each $C^f_\gamma(\tau_0)$ depends on an auxiliary
basepoint $\tau_0$ and on a group element $\gamma$, while the period
polynomial must refer to neither.  Degree-one group cohomology is the
standard bookkeeping for exactly this situation, and only three
definitions are needed.  A \emph{$1$-cocycle} is a function
$c\colon \Gamma \to V_n$ satisfying $c_{gh} = c_g|_h + c_h$, and
$\gamma \mapsto C^f_\gamma(\tau_0)$ is one for each fixed $\tau_0$
(Lemma~\ref{lem:cocycle}).  A \emph{$1$-coboundary} is the special
cocycle $\delta P\colon \gamma \mapsto P|_\gamma - P$ built from a
polynomial $P \in V_n$; by \eqref{eq:bp-cocycle-coboundary}, shifting
$\tau_0 \to \tau_0'$ (no pole crossed) shifts the whole family
$\{C^f_\gamma\}_\gamma$ by one such $\delta P$ --- the same $P$ for
every $\gamma$, which is why the quotient retains anything at all.  The
\emph{cohomology class}
$[C^f] \in H^1(\Gamma; V_n) :=
\{\text{cocycles}\}/\{\text{coboundaries}\}$ is therefore
basepoint-free, and a \emph{representative} of it is any cocycle
$C^f - \delta P$.  This kills the
$\tau_0$. %CLAUDE: replaces the former lines 313--320, which described coboundaries as polynomials and the equivalence as acting on each C^f_\gamma separately.

A reader outside the field might object that no machinery should be
needed: the classical period polynomial is a single integral along one
$S$-contour, so why not take the finite $S$-contour $C^f_S(\tau_0)$ and
stop?  Because at a finite basepoint the raw $S$-value is pure
redundancy.  Torsion alone forces
$C^f_S\big|(1+S) = C^f_{S^2} = C^f_{-I} = 0$ at \emph{every} $\tau_0$,
and since $S$ acts on $V_n$ as an involution, its $(-1)$-eigenspace is
also the image of $\delta_S$: the value $C^f_S$ lies in
$\ker(1+S) = \im \delta_S$, exactly the space of shifts that basepoint
motion induces on it.  (At $k = 12$ this kernel is six-dimensional
inside the eleven-dimensional $V_{10}$; membership selects a subspace,
never a point.)  The two relations cutting out the Kohnen--Zagier space
$W_{k-2}$ \cite{KZ1984} therefore have different statuses: the
$S$-relation is \emph{free}, holding before any choice is made, while
the $U$-relation must be \emph{purchased} --- it appears only after the
representative is pinned, and pinning is the role of the other
generator.

The generator $T$ is built for the job.  It has infinite order and acts
on $V_n$ unipotently --- a single Jordan block --- so
$\delta_T := (\,\cdot\,)|_T - \mathrm{id}$ has one-dimensional kernel
and cokernel: the equation has one scalar of obstruction and one scalar
of slack.  We choose, within the class, the representative
$C^f - \delta P$ \emph{whose $T$-value vanishes}, i.e.\ solve
$\delta_T P = (2\pi i)^{-(n+1)} C^f_T(\tau_0)$ on $V_n$.  Solvability is
itself arithmetic: unit-segment $T$-periodicity identifies the
obstruction in $V_n / \im \delta_T \cong \CC \cdot [X^n]$ as the
constant Fourier mode (Lemma~\ref{lem:solvability-cuspidality}),
$[C^f_T]_{X^n} = (2\pi i)^{n+1} c_f(0)$ --- the residue of $\omega_f$ at
the very cusp that $T$ stabilizes.  The normalized representative exists
if and only if $c_f(0) = 0$; classes that pass are the \emph{parabolic}
ones, identified by Eichler--Shimura \cite{KZ1984} with
$S_k \oplus \overline{S_k}$, the lone obstructed direction being
Eisenstein.  Cuspidality, in this framework, is not a convergence
hypothesis; it is the solvability condition for the normalization.

With the $T$-value normalized away, the entire cuspidal content sits on
the surviving generator value,
\begin{equation}\label{eq:rf-normalized}
  r_f \;:=\; C^f_S(\tau_0) \;-\; (2\pi i)^{n+1}\bigl(P|_S - P\bigr),
\end{equation}
well-defined up to the line $\CC\,(X^n - Y^n)$ left over from
$\ker \delta_T = \CC\, Y^n$ --- an ambiguity touching only the boundary
monomials, which is why Theorem~\ref{thm:periods} ranges over the
interior $\ell$ alone.  The $U$-relation now arrives, purchased as
promised: conditioned on the vanishing $T$-value, the group relation
$(ST)^3 = -I$ becomes $r_f|\bigl(1 + (TS) + (TS)^2\bigr) = 0$, and $r_f$
\emph{is} the period polynomial.  Indeed, for holomorphic $f \in S_k$,
sending $\tau_0 \to i\infty$ kills $C^f_T$ outright and
\eqref{eq:rf-normalized} collapses to the classical integral
\eqref{eq:classical-rf} --- whose familiar freedom from $\tau_0$ and
$\gamma$ is thereby revealed as the silent choices $\gamma = S$,
$\tau_0 = i\infty$.  Its interior coefficients are the periods against
Brown's rational basis,
$[r_f]_{X^{k-2-\ell}Y^\ell} = \omega^\pm(f)\, c_S^{\,\ell}$; the precise
coordinate statement is given in Section~\ref{sec:cocycle}.  The finite
$S$-contour was the right object all along, but it requires the
correction term in \eqref{eq:rf-normalized}, sourced entirely by the
$T$-contour --- and producing that correction in closed form is the
business of this paper.

% ==== END PASTE: BLOCK 1 ====

\blockrule
\begin{whereitgoes}
The existing line 353, ``The coboundary equation admits a closed-form
inverse --- the polynomial inversion of [DR2017]\ldots'', follows the
final sentence above directly; ``that correction in closed form'' is its
antecedent.  For meromorphic $f$ with interior poles, the identity
$C^f_S|(1+S) = 0$ holds per homotopy component; state that precisely in
Section~\texttt{sec:wall-crossing}, not here.
\end{whereitgoes}

% ============================================================================
\section*{Block 2 --- replaces lines 767--774 (Section \texttt{sec:cocycle})}
% ============================================================================
\blockrule

% ==== BEGIN PASTE: BLOCK 2 (replaces lines 767--774; delete the %DAM note at line 774) ====

\paragraph{Parabolic cohomology and periods.}
The cohomology class $[C^f] \in H^1(\Gamma; V_n)$ is independent of
$\tau_0$ (on the connected component) and of the choice of $\delta P$
representative.  Its restriction to the cusp stabilizer
$\langle T \rangle \cong \ZZ$ lands in
$H^1(\langle T \rangle; V_n) = V_n / \im \delta_T \cong \CC\cdot[X^n]$,
and by Lemma~\ref{lem:solvability-cuspidality}(c) the restricted class is
$(2\pi i)^{n+1}\, c_f(0) \cdot [X^n]$: the residue of $\omega_f$ at the
cusp, valued in the $T$-coinvariants.  The class is \emph{parabolic} ---
cohomologous to a cocycle with vanishing $T$-value --- if and only if
this residue vanishes, i.e.\ if and only if $f \in S^!_k$.  The classical
Eichler--Shimura theorem \cite{KZ1984} identifies
$H^1(\Gamma; V_n) \cong S_k \oplus \overline{S_k} \oplus
\CC_{\mathrm{Eis}}$ and the parabolic subspace
$\Hpar(\Gamma; V_n) \cong S_k \oplus \overline{S_k}$; on cocycles
normalized so that $c_T = 0$, the assignment $c \mapsto c_S$ further
identifies $\Hpar(\Gamma; V_n) \cong W_{k-2}/\CC\,(X^n - Y^n)$, with the
quotient line accounting for the residual freedom
$P \mapsto P + \CC\, Y^n$.  When $\dim S_k = 1$, each parity eigenspace
of $\Hpar$ is one-dimensional, and the two coordinates of a parabolic
class in Brown's rational basis $p^\pm_f$ are its periods:
$[C^f]^\pm = \omega^\pm(f)[p^\pm_f]$ for $f \in S_k$, while the
quasi-periods $\eta^\pm$ arise as the coordinates of the class
$[C^{f'}]$ of a weakly holomorphic Hecke partner $f' \in S^!_k$ (at
$k = 12$: $f = \Delta$, $f' = \Dhat$)
\cite{Brown2017}. %CLAUDE: corrected -- the previous version presented \omega^\pm and \eta^\pm jointly as "the linear coordinates of [C^f]"; a single class in the two-dimensional space \Hpar has two coordinates, and the four numbers of the period matrix come from the two classes [C^\Delta] and [C^\Dhat].

% ==== END PASTE: BLOCK 2 ====

\blockrule

% ============================================================================
\section*{Block 3 --- sign fix in \texttt{eq:bp-cocycle-coboundary}
(lines 755--760)}
% ============================================================================

\noindent
Minimal fix: in the final clause of the display, swap the arguments of
$F$, replacing $P = (2\pi i)^{n+1}\bigl(F(\tau_0') - F(\tau_0)\bigr)$ by
\[
  P \;=\; (2\pi i)^{n+1}\bigl(F(\tau_0) - F(\tau_0')\bigr) \in V_n
  \otimes \CC .
\]
Derivation (two lines, via the equivariance step of
Lemma~\ref{lem:cocycle}, $\int_{\gamma^{-1}a}^{\gamma^{-1}b}\omega_f =
\bigl(\int_a^b \omega_f\bigr)\big|_\gamma$): splitting
$\int_{\gamma^{-1}\tau_0'}^{\tau_0'} =
\int_{\gamma^{-1}\tau_0'}^{\gamma^{-1}\tau_0} +
\int_{\gamma^{-1}\tau_0}^{\tau_0} + \int_{\tau_0}^{\tau_0'}$ gives
\[
  C^f_\gamma(\tau_0') - C^f_\gamma(\tau_0)
  \;=\; (2\pi i)^{n+1}\Bigl[\Bigl(\textstyle\int_{\tau_0'}^{\tau_0}
  \omega_f\Bigr)\Big|_\gamma + \int_{\tau_0}^{\tau_0'}\omega_f\Bigr]
  \;=\; A - A|_\gamma,
  \qquad
  A := (2\pi i)^{n+1}\!\!\int_{\tau_0}^{\tau_0'}\!\omega_f ,
\]
which is $P|_\gamma - P$ for $P = -A$.  The cohomological conclusion is
unaffected; only the labelling of $P$ changes.  Check the convention once
at the single downstream point where $P$ appears explicitly, the
normalized value \eqref{eq:rf-normalized}.

\end{document}
